% Source: https://artofproblemsolving.com/community/c3h2416978_1200th_post_dedicated_to_bloggy
% Stable identity: aops-topic-2416978-q15
Replacing $x$ by $1-x$ gives a second equation:
\[f(x)+xf(1-x)=x,\qquad f(1-x)+(1-x)f(x)=1-x.\]
From the second, $f(1-x)=(1-x)-(1-x)f(x)$. Substituting into the first,
\[f(x)+x(1-x)-x(1-x)f(x)=x\quad\Longrightarrow\quad\left(1-x+x^2\right)f(x)=x-x(1-x)=x^2.\]
Since $x^2-x+1=\left(x-\dfrac12\right)^2+\dfrac34>0$, we get
\[\boxed{f(x)=\dfrac{x^2}{x^2-x+1}}.\]
This function works: $(1-x)^2-(1-x)+1=x^2-x+1$, so $f(1-x)=\dfrac{(1-x)^2}{x^2-x+1}$ and
\[f(x)+xf(1-x)=\dfrac{x^2+x(1-x)^2}{x^2-x+1}=\dfrac{x\left(x^2-x+1\right)}{x^2-x+1}=x.\]
So it is the unique solution.
